\Needspace{.85\textheight}
\section*{Разобранный пример: сумма первых натуральных чисел}

\begin{theorem}[Сумма первых натуральных чисел]
Для каждого \(n\in\mathbb N\) выполнено
\[
1+2+\dots+n=\frac{n(n+1)}2.
\]
При \(n=1\) слева находится единственное слагаемое \(1\).
\end{theorem}

\begin{proof}
Обозначим доказываемое равенство через \(P(n)\).

\textbf{1. База.} При \(n=1\) левая часть равна \(1\), а правая ---
\(\frac{1\cdot2}{2}=1\). Значит, \(P(1)\) верно.

\textbf{2. Предположение.} Зафиксируем произвольное \(k\ge1\) и предположим
\[
1+2+\dots+k=\frac{k(k+1)}2.
\]
Это временное предположение только для одного произвольного \(k\).

\textbf{3. Переход.} Сумма первых \(k+1\) чисел получается добавлением
одного нового слагаемого \(k+1\). Поэтому
\[
\begin{aligned}
1+2+\dots+k+(k+1)
&=\frac{k(k+1)}2+(k+1)\\
&=\frac{k(k+1)+2(k+1)}2\\
&=\frac{(k+1)(k+2)}2.
\end{aligned}
\]
В первой строке использовано индукционное предположение; далее приведены к общему
знаменателю два слагаемых и вынесен общий множитель \(k+1\).
Последняя строка --- в точности формула \(P(k+1)\), потому что
\((k+1)+1=k+2\). Общий переход доказан.

База и переход вместе дают по методу математической индукции требуемое равенство
для каждого \(n\ge1\).
\end{proof}

Индукция объясняет, почему формула верна для всех натуральных \(n\).
Подстановка, например, \(n=5\) даёт \(1+2+3+4+5=15\), но сама по себе
проверяет лишь один случай и не заменяет проведённого доказательства.
